% !TeX root = ../papers/nc-p3.tex
\section{純関係を持つ辺の分布と長さ二以上の障害}
\label{sec:sparse-pure-edge-exclusion}
\begin{theorem}\label{thm:sparse-pure-edge-exclusion}
標準正則四窓で純関係が一つでもあれば、純関係を持つ辺は少なくとも二つあり、
そのうち一辺の純関係スキームは長さ2以上である。従って四辺合計の長さ2以下は
空の場合に限り、重複点だけを持つ一辺も任意の長さで排除される。
\end{theorem}
\begin{proof}
定理\ref{thm:finite-pure-relation-locus}により有限性は既知で、
曲面・有限辺・一点グラフの排除は定理\ref{thm:single-pure-branch-exclusion}の冒頭による。
同定理の列挙を、全辺の純関係スキームが長さ1以下の場合と、一辺だけ任意の長さを許す場合に拡張する。
中心の例外固有方向の重複込み個数を $l=3-e$ と置く。
一中心では $0\le l\le3$、同位相二中心では各 $l\le2$ である。
その和は当該辺の純関係スキームの長さを超えない。一辺にだけ純関係を許すときは、
その辺に対する上界4で全ての中心型を列挙できる。これは辺自体の長さの上界を仮定するものではない。
全辺長さ1以下の場合、記述の位相と向きを選んで位置0の最初の順中心を例外にできる。

次数列の整数条件の変更点を記す。$l>0$ の中心に対し、同位相の初回以外の移動は
$e>0$ のときには主曲線（次数 $e$）に一回以下、または例外直線（それぞれ次数1）に合計 $l$ 回以下である。
$e=1$ なら主曲線と例外直線を合わせて高々三回の直線への移動を数える。
$e=0$ では非定数の主曲線がない。初回直線の外にある次位相の順中心も高々 $l$ 個であり、
同位相が二中心のときは引き続き第二次数が相手の $e$ に等しい。
普通中心がどちらかの向きにあれば $m\le12$、全てが例外なら
単一辺だけの設定では二つ以上の占有位相を得られず、全辺単純の設定では
一位相の移動が高々三回なので $m\le16$ である。

非初回の移動の始点曲線を $D$、その一つ前の曲線が埋込みであるとする。
移動先が主曲線なら初回直線 $L$ からの射影が双有理であり、
定理\ref{thm:branch-cubic-contact}の引戻しの議論で $\deg D\le2$ となる。
例外方向にあると分かる場合は $D$ が $L$ を含む平面に入るから、
$D\ne L$ より $D$ 上の $L$ の固定因子の長さ自体が $\deg D$ となり、同じ議論で $\deg D\le1$ である。
主曲線も直線なら緩い方の上界2を使えばよい。
ここで埋込みと数えるのは次数1,2の全て、初回から二番目と終回から二番目、
および反対三次曲線の像である。その他の三次曲線に埋込みを仮定しない。

以上と前定理の終点対応、次数和、三階差の符号を使う。単純な一純関係を両向きの例外中心が共有すれば
その二中心を二直線列で結ぶ。連続二直線列の排除は過剰接触を生じる中央辺の長さが1以下の場合に
両向きで適用する。同位相の二つの $(1,3,1)$ が一つの反対中心を通れない条件も両向きで使う。
添付コードは各中心に許される $l$ を全て選び、上の上限内の全列を相異なる終点に割り当てる。
結果は両向きとも相対二位相に一中心ずつで、次の候補に限られる。
一辺だけを許した場合は両方 $(1,3,1)$ であり、前定理の最後の証明により不可能である。
全辺長さ1以下の場合もこれと同じか、一方 $(1,3,1)$、他方
$(1,3,a,m,c,3,1)$ である。後者では $1\le a,m,c\le2$、
長い列の始点の順鉛筆と終点の逆鉛筆はともに $L_3$ 型である。

長い列の第二・第六三次曲線は同じ位相で相異なり、互いに相手を初回に持つ向きの反対三次曲線となる。
第三曲線と第五曲線は相手側の初回直線ではないので、これは二つの $L_3$ の二者択一に従う。
定理\ref{thm:branch-cubic-contact}の単元を用いた小行列式の包含により、
両者の交叉はそれぞれ逆・順の中心に欠損の長さで含まれる部分因子を持つ。
その長さの和は少なくとも $(m-3a+8)+(m-3c+8)\ge6$ である。
二中心は異なるので、相異なるねじれ三次曲線の交叉長の上界5に反する。
\end{proof}
